\section{Contribution, limitations and further questions}

The direct mixed-product proof makes the cancellation mechanism explicit and gives an optimal universal prefactor for the stated general class. The diagonal-moment results separate three distinct questions: when the published leading term is accurate, what replaces it at the transition, and how far that replacement can remain accurate.

The base recurrence, its normalized parity form and the leading coefficients belong to Collins and Nagatsu \cite{CN}. The mixed estimate itself is already part of Bordenave's spectral-gap method \cite{Bordenave}. The special-function identity behind the positive partition expansion is also classical. The contributions established here are the direct quantitative proof, the necessity and bounded-comparability characterization of the cubic scale, the wider relative hyperbolic window, and its square-root correction. A targeted literature check found no exact duplicate of these quantitative refinements, but cannot certify the absence of every earlier equivalent asymptotic formula.

Bordenave and Zhu \cite[Lemma 2.1]{BZ} retain the original singleton factor $3k/\sqrt N$ and explicitly bound the raw factorial ratio by $\sqrt e$ in their equation (2.3). That elementary bound is a known ingredient here. Our mixed-product refinement combines the exact correction $F_n(a-h)$ with the smaller singleton factor $\beta$, independently of repeated multiplicities, and proves optimality of the resulting universal prefactor for the stated class.

The Laguerre representation makes the parameter regime precise: the degree is $k$, the parameter is $\alpha=-N-1$, and the argument is $-N$. Thus $-\alpha/k\to\infty$ on our growing-degree square-root ranges. The negative-parameter regime described by Dai and Wong \cite{DW} has $-\alpha/k$ tending to a finite limit greater than one. Theorem 1 of Dai, Ismail and Wang \cite{DIW} instead assumes $\alpha>0$ with $\alpha/k\to\infty$. These statements do not directly supply the present uniform approximation. A polynomial identity or analytic continuation alone does not transfer their uniform error estimates to our regime.

The inequalities are proved symbolically for their stated parameter ranges. Exact finite computations check normalization, the depleted recurrence and the conditioning examples. Numerical plots illustrate the asymptotic regimes. Neither kind of computation is used as a substitute for a general proof.

The remaining graph-level task is to combine the sharper cell-pattern estimate with circuit enumeration and tangle remainders while retaining the degree dependence. A better estimate for one fixed product does not automatically improve a spectral theorem: the patterns that dominate the full sum must be identified first. The present paper does not assert a new second-eigenvalue bound or an end-to-end audit of either random-graph proof.
